\scanpage{396}\mspage{310}

$\pi'_0$-accessible et que $\pi>\pi'_0$ satisfait $\unsure{\pi_0|\pi}^{\pi'_0}=\pi$, on en conclut que $\underline S$ est $\underline S_\pi$-accessible, i.e. $\underline S\xleftarrow{\approx}\mathrm{Ind}_\pi(\underline S_\pi)$. Donc en vertu de 4.13.3, il reste à voir que $\underline S_\pi\subset M_\pi$.

\deleted{Or \inserted{comme $\pi=\pi^{\pi'_0}$,} $\underline S_\pi$ est formé des \inserted{objets de $\underline S$ \unsure{isom.} à des} $\underset{\longrightarrow}{\lim}$ $x_i$, où $I$ \unsure{qui}} ou ce qui revient au même, que $\underline S_\pi\subset\underline S(\pi)\overset{\text{\unsure{def|déf}}}{=}\underline S\cap M_\pi$. (En fait, on aura $=$ égalité, car il est \deleted{\unsure{immédiat}} \inserted{clair} que $\underline S(\pi)\subset\underline S_\pi\;\ldots$).

Or $\underline S(\pi)$ est une sous-catégorie de $\underline S$ stable par $\underset{\longrightarrow}{\lim}$ grdes devant $\pi'_0$ et de cardinal $\leq\pi$ (puisque $\underline S$ et $M_\pi$ sont stables par ces limites) \uline{dans $M$}, et que ces limites dans $\underline S$ et dans \uline{$M$} sont \unsure{itou|idem}), et $\underline S(\pi)$ contient $\underline S_{\pi'_0}$. Donc $\underline S(\pi)$ contient $\underline S_\pi$, en vertu de 4.7.1 (page 161), implication c)$\Longrightarrow$b', appliqué à $\underline S,\pi'_0,\pi$ \unsure{ald|a.l.d.} $M,\pi_0,\pi$. Cela prouve \unsure{bien|donc} le corollaire 4.13.16.

\uline{Démonstration de} 4.13.15. Voici les implications $\pm$ tautologiques

\diagram{\begin{tikzcd}[ampersand replacement=\&,column sep=small,row sep=0.4em]
\& \& \& \& \deleted{\unsure{d}}c'\\
(*)\qquad \& a)\arrow[r,Rightarrow] \& \deleted{\unsure{b}}a'\arrow[r,Rightarrow,"\unsure{\cdot}"] \& \deleted{\unsure{d}}c\arrow[ru,Rightarrow]\arrow[rd,Rightarrow] \& \\
\& \& \& \& b,
\end{tikzcd}}

% archive sheet number: 395
